\Needspace{11\baselineskip}
\section{Lebesgue--Nikodym 定理的一个证明}\label{sec:yosida-03-08}
\noindent\textit{原书位置：中译本第 97--99 页；英文版第 93--95 页。}

\begin{theorem}[非负测度的 Radon--Nikodym 密度]\label{thm:y3-radon-nikodym}
设 $\mu,\nu$ 为同一可测空间 $(S,\Sigma)$ 上的 $\sigma$ 有限非负测度，且 $\nu\ll\mu$，即 $\mu(E)=0$ 蕴含 $\nu(E)=0$。则存在非负可测函数 $p$，在 $\mu$ 几乎处处意义下唯一，使
\[
 \nu(E)=\int_Ep\,d\mu\quad(E\in\Sigma).
\]
等式允许两端为无穷。原书以 Lebesgue--Nikodym 为名，本稿沿用节名并注明现今常用名称。
\end{theorem}
\begin{proof}
先设两测度均有限，置 $\rho=\mu+\nu$，使用实 Hilbert 空间 $L^2(\rho)$。泛函 $L(f)=\int f\,d\nu$ 良定义且有界，因为
\[
 |L(f)|\le\nu(S)^{1/2}\|f\|_{L^2(\nu)}
 \le\nu(S)^{1/2}\|f\|_{L^2(\rho)}.
\]
Riesz 定理给出实函数 $h\in L^2(\rho)$，使 $\int f\,d\nu=\int fh\,d\rho$。对特征函数测试，得 $\nu(E)=\int_Eh\,d\rho$。由于 $0\le\nu(E)\le\rho(E)$ 对每个 $E$ 成立，在集合 $\{h<0\}$ 与 $\{h>1\}$ 上测试（或先取阈值子集），得 $0\le h\le1$，$\rho$ 几乎处处。相减还得 $\mu(E)=\int_E(1-h)\,d\rho$。

在 $A=\{h=1\}$ 上 $\mu(A)=0$，绝对连续性给出 $\nu(A)=0$，故 $\rho(A)=0$。修改零测集后可设 $0\le h<1$，令 $p=h/(1-h)$。由对特征函数、简单函数和单调极限依次成立的积分换测度公式，
\[
 \int_Ep\,d\mu=\int_Ep(1-h)\,d\rho
 =\int_Eh\,d\rho=\nu(E).
\]
这一步只使用简单函数逼近与单调收敛，并没有预先调用待证密度定理。

一般情形中，$\mu$ 与 $\nu$ 的有限测度覆盖取两两交，再枚举并去除此前各集，得到互不相交的可测分片 $E_j$ 覆盖 $S$，每片的两种测度均有限。在每片应用有限情形得到 $p_j$，拼成 $p=\sum_j1_{E_j}p_j$；可数可加性及非负积分的可数可加性给出全空间等式。

为证唯一，若 $p,q$ 都给出表示，先限制到一个 $\nu$ 有限的覆盖片，并测试其内 $\{p\ge q+1/r\}$。两密度在此可积且积分相同，故该集合的 $\mu$ 测度为零。对 $r$ 及覆盖片取可数并，再交换 $p,q$，即得 $p=q$ 几乎处处。密度在 $\nu$ 有限片上可积，亦保证它 $\mu$ 几乎处处有限。
\end{proof}

图 \ref{fig:y3-radon-nikodym} 只针对有限测度步骤，帮助追踪 $h$、$1-h$ 与最终密度 $p$ 的关系。
\begin{figure}[htbp]
\centering
\begin{tikzpicture}[fa diagram]
  \node[fa box] (rho) at (0,2.3) {$\rho=\mu+\nu$};
  \node[fa box] (mu) at (-3.6,0) {$\mu(E)=\int_E(1-h)\,d\rho$};
  \node[fa box] (nu) at (3.6,0) {$\nu(E)=\int_Eh\,d\rho$};
  \draw[->,fa main] (rho)--(mu) node[midway,left] {$1-h$};
  \draw[->,fa main] (rho)--(nu) node[midway,right] {$h$};
  \draw[->,fa accent] (mu.south) to[bend right=22] node[midway,below] {$p=h/(1-h)$} (nu.south);
  \node[fa note] at (0,3.5) {有限测度情形：Riesz 表示先得到 $h$，$0\le h\le1$};
  \node[fa note] at (0,-2.5) {$A=\{h=1\}$：$\mu(A)=0\Rightarrow\nu(A)=0\Rightarrow\rho(A)=0$。};
\end{tikzpicture}
\caption{Radon--Nikodym 证明的共同参照测度。箭头上的函数表示积分权重；绝对连续性排除分母为零的集合，随后用简单函数与单调收敛证明换测度公式。一般 $\sigma$ 有限情形先分成两测度都有限的片。}
\label{fig:y3-radon-nikodym}
\end{figure}



\begin{example}[原子与绝对连续性]
在可数集上，若 $\mu(\{j\})=a_j>0$，$\nu(\{j\})=b_j\ge0$ 且各值有限，则 $p(j)=b_j/a_j$。相反，$\R$ 上的 Dirac 测度 $\delta_0$ 不绝对连续于 Lebesgue 测度，因为单点的后者测度为零而前者为一；不存在相应的 Lebesgue 可积密度。
\end{example}

\begin{exercise}
设 $d\nu=p\,d\mu$，其中 $p\ge0$。证明对每个非负可测 $g$ 有 $\int g\,d\nu=\int gp\,d\mu$。若又有 $d\lambda=q\,d\nu$，推出 $d\lambda=qp\,d\mu$。
\end{exercise}
\begin{solution}
对特征函数公式是密度定义，对非负简单函数由有限线性组合成立。取简单函数 $g_n\uparrow g$，两边应用单调收敛即可；零密度处乘积约定为零。再把 $g=1_Eq$ 代入，得 $\lambda(E)=\int_Eq\,d\nu=\int_Eqp\,d\mu$，这就是密度的链式法则。
\end{solution}
